% Options for packages loaded elsewhere
\PassOptionsToPackage{unicode}{hyperref}
\PassOptionsToPackage{hyphens}{url}
%
\documentclass[
]{article}
\usepackage{amsmath,amssymb}
\usepackage{lmodern}
\usepackage{iftex}
\ifPDFTeX
  \usepackage[T1]{fontenc}
  \usepackage[utf8]{inputenc}
  \usepackage{textcomp} % provide euro and other symbols
\else % if luatex or xetex
  \usepackage{unicode-math}
  \defaultfontfeatures{Scale=MatchLowercase}
  \defaultfontfeatures[\rmfamily]{Ligatures=TeX,Scale=1}
\fi
% Use upquote if available, for straight quotes in verbatim environments
\IfFileExists{upquote.sty}{\usepackage{upquote}}{}
\IfFileExists{microtype.sty}{% use microtype if available
  \usepackage[]{microtype}
  \UseMicrotypeSet[protrusion]{basicmath} % disable protrusion for tt fonts
}{}
\makeatletter
\@ifundefined{KOMAClassName}{% if non-KOMA class
  \IfFileExists{parskip.sty}{%
    \usepackage{parskip}
  }{% else
    \setlength{\parindent}{0pt}
    \setlength{\parskip}{6pt plus 2pt minus 1pt}}
}{% if KOMA class
  \KOMAoptions{parskip=half}}
\makeatother
\usepackage{xcolor}
\setlength{\emergencystretch}{3em} % prevent overfull lines
\providecommand{\tightlist}{%
  \setlength{\itemsep}{0pt}\setlength{\parskip}{0pt}}
\setcounter{secnumdepth}{-\maxdimen} % remove section numbering
\ifLuaTeX
  \usepackage{selnolig}  % disable illegal ligatures
\fi
\IfFileExists{bookmark.sty}{\usepackage{bookmark}}{\usepackage{hyperref}}
\IfFileExists{xurl.sty}{\usepackage{xurl}}{} % add URL line breaks if available
\urlstyle{same} % disable monospaced font for URLs
\hypersetup{
  hidelinks,
  pdfcreator={LaTeX via pandoc}}

\author{}
\date{}

\begin{document}

\documentclass{article}
\usepackage{amsmath}
\usepackage{amsfonts}
\begin{document}

\section*{Lecture Scribe}

\subsection*{CSE400 – Fundamentals of Probability in Computing}

\subsection*{Lecture 9: Uniform, Exponential, Laplace and Gamma Random Variables}

---

\section*{Outline}

The lecture covers the following topics:

\subsubsection*{Types of Continuous Random Variables}

\begin{itemize}
\item Uniform Random Variable: Example
\item Exponential Random Variable: Example
\item Laplace Random Variable: Example
\item Gamma Random Variable
\begin{itemize}
\item Graph and Special Cases
\item Example
\item Homework Problem
\end{itemize}
\end{itemize}

\subsubsection*{Problem Solving}

---

\section*{1. Types of Continuous Random Variables}

\subsection*{1.1 Uniform Random Variable}

The uniform random variable is defined by the following PDF:

\[
f_X(x)=
\begin{cases}
\dfrac{1}{b-a}, & a \le x < b,[4pt]\\
0, & \text{elsewhere}
\end{cases}
\]

Its CDF is

\[
F_X(x)=
\begin{cases}
0, & x<a,[4pt]\\
\dfrac{x-a}{b-a}, & a \le x < b,[6pt]\\
1, & x \ge b.
\end{cases}
\]

The slides also present the graph of the PDF and the graph of the CDF for the uniform random variable.

---

\section*{2. Example \#1 (Uniform RV)}

\subsection*{Problem}

The phase of a sinusoid, $(\Theta)$, is uniformly distributed over $([0,2\pi))$ so that its PDF is

\[
f_\Theta(\theta)=
\begin{cases}
\dfrac{1}{2\pi}, & 0 \le \theta < 2\pi,[4pt]\\
0, & \text{otherwise.}
\end{cases}
\]

Find:

\begin{itemize}
\item (a) $\left(\Pr!\left(\theta>\dfrac{3\pi}{4}\right)\right)$
\item (b) $\left(\Pr!\left(\theta<\pi \mid \theta>\dfrac{3\pi}{4}\right)\right)$
\item (c) $\left(\Pr!\left(\cos(\theta)<\dfrac{1}{2}\right)\right)$
\end{itemize}

\subsection*{Solution}

\subsubsection*{Given}

\[
f_\Theta(\theta)=
\begin{cases}
\dfrac{1}{2\pi}, & 0 \le \theta < 2\pi,[4pt]\\
0, & \text{otherwise}
\end{cases}
\]

For a uniform RV on $([0,2\pi))$,

\[
\Pr(a<\Theta<b)=\frac{b-a}{2\pi}
\]

\subsubsection*{(a) Compute $\left(\Pr!\left(\Theta>\dfrac{3\pi}{4}\right)\right)$}

\[
\Pr!\left(\Theta>\frac{3\pi}{4}\right)
\]
\[
======================================
\]
\[
\# \frac{2\pi-\frac{3\pi}{4}}{2\pi}
\]
\[
\frac{5}{8}
\]

\subsubsection*{(b) Compute $\left(\Pr!\left(\Theta<\pi \mid \Theta>\dfrac{3\pi}{4}\right)\right)$}

Using

\[
\Pr(A\mid B)=\frac{\Pr(A\cap B)}{\Pr(B)}
\]

the numerator is

\[
\Pr!\left(\frac{3\pi}{4}<\Theta<\pi\right)
\]
\[
==========================================
\]
\[
\# \frac{\pi-\frac{3\pi}{4}}{2\pi}
\]
\[
\frac{1}{8}
\]

and

\[
\Pr(B)=\Pr!\left(\Theta>\frac{3\pi}{4}\right)=\frac{5}{8}
\]

Hence,

\[
\Pr!\left(\Theta<\pi \mid \Theta>\frac{3\pi}{4}\right)
\]
\[
======================================================
\]
\[
\# \frac{1/8}{5/8}
\]
\[
\frac{1}{5}
\]

\subsubsection*{(c) Compute $\left(\Pr!\left(\cos(\Theta)<\dfrac{1}{2}\right)\right)$}

\[
\cos\Theta=\frac{1}{2}
\Rightarrow
\Theta=\frac{\pi}{3},\ \frac{5\pi}{3}
\]

Thus,

\[
\cos\Theta<\frac{1}{2}
\quad \text{for} \quad
\frac{\pi}{3}<\Theta<\frac{5\pi}{3}
\]

Therefore,

\[
\Pr!\left(\cos(\Theta)<\frac{1}{2}\right)
\]
\[
=========================================
\]
\[
\# \frac{\frac{5\pi}{3}-\frac{\pi}{3}}{2\pi}
\]
\[
\# \frac{4\pi/3}{2\pi}
\]
\[
\frac{2}{3}
\]

---

\section*{3. Uniform Random Variable: Application Examples}

The slides list the following application examples:

\begin{itemize}
\item The phase of a sinusoidal signal, when all phase angles between $(0)$ and $(2\pi)$ are equally likely.
\item A random number generated by a computer between 0 and 1 for simulations.
\item The arrival time of a user within a known time window, assuming no time preference.
\end{itemize}

---

\section*{4. Exponential Random Variable}

For any $(b>0)$, the exponential random variable has PDF and CDF given by

\[
f_X(x)=\frac{1}{b}\exp!\left(-\frac{x}{b}\right)u(x)
\]

and

\[
F_X(x)=\left[1-\exp!\left(-\frac{x}{b}\right)\right]u(x)
\]

The slides also present the graph of the PDF and the graph of the CDF for the exponential random variable, with the CDF graph labeled $((b=2))$.

---

\section*{5. Example \#2 (Exponential RV)}

\subsection*{Problem}

Let $(X)$ be an exponential random variable with PDF

\[
f_X(x)=e^{-x}u(x)
\]

Find:

\begin{itemize}
\item (a) $(\Pr(3X<5))$
\item (b) Generalize your answer to part (a) to find $(\Pr(3X<y))$ for some arbitrary constant $(y)$.
\item (c) Note that if we define a new random variable according to $(Y=3X)$ then your answer to part (b) is the CDF of $(Y)$, $(F_Y(y))$. Given your answer to part (b), find $(f_Y(y))$.
\end{itemize}

\subsection*{Solution}

\subsubsection*{Given}

\[
f_X(x)=e^{-x}u(x)
\Rightarrow
F_X(x)=\left[1-e^{-x}\right]u(x)
\]

\subsubsection*{(a) Compute $(\Pr(3X<5))$}

\[
\Pr(3X<5)
\]
\[
=========
\]
\[
\# \Pr!\left(X<\frac{5}{3}\right)
\]
\[
\# F_X!\left(\frac{5}{3}\right)
\]
\[
1-e^{-5/3}
\]

\subsubsection*{(b) Compute $(\Pr(3X<y))$}

\[
\Pr(3X<y)
\]
\[
=========
\]
\[
\Pr!\left(X<\frac{y}{3}\right)
\]

\[
\# [
\]

\[
\# F_X!\left(\frac{y}{3}\right)
\]
\[
\left[1-e^{-y/3}\right]u(y)
\]

\subsubsection*{(c) Let $(Y=3X)$. Find $(f_Y(y))$}

From part (b), the CDF of $(Y)$ is

\[
F_Y(y)=\left[1-e^{-y/3}\right]u(y)
\]

Differentiating,

\[
f_Y(y)=\frac{d}{dy}F_Y(y)
\]
\[
=========================
\]
\[
\frac{1}{3}e^{-y/3}u(y)
\]

Hence, $(Y)$ is an exponential random variable with parameter

\[
\lambda=\frac{1}{3}
\]

---

\section*{6. Exponential Random Variable: Application Examples}

The slides list the following application examples:

\begin{itemize}
\item Time until a radioactive particle decays, assuming a constant decay rate.
\item Duration of an idle period of a wireless communication channel before it becomes busy.
\item Time until the next earthquake in a simplified seismic activity model with a constant occurrence rate.
\end{itemize}

---

\section*{7. Laplace Random Variable}

The Laplace random variable is given by the following PDF:

\[
f_X(x)=\frac{1}{2b}\exp!\left(-\frac{|x|}{b}\right)
\]

Its CDF is

\[
F_X(x)=
\begin{cases}
\dfrac{1}{2}\exp!\left(\dfrac{x}{b}\right), & x<0,[8pt]\\
1-\dfrac{1}{2}\exp!\left(-\dfrac{x}{b}\right), & x\ge 0
\end{cases}
\]

The slides state the following application:

The Laplace random variable has been used to model the amplitude of a speech (voice) signal.

The slides also present the graph of the PDF and the graph of the CDF for the Laplace random variable, with the CDF graph labeled $((b=2))$.

---

\section*{8. Example \#3 (Laplace RV)}

\subsection*{Problem}

Let $(W)$ be a Laplace random variable with a PDF given by

\[
f_W(w)=ce^{-2|w|}
\]

Find:

\begin{itemize}
\item (a) the value of the constant $(c)$
\item (b) $(\Pr(-1<W<2))$
\item (c) $(\Pr(W>0\mid -1<W<2))$
\end{itemize}

\subsection*{Solution}

\subsubsection*{Given}

\[
f_W(w)=ce^{-2|w|}
\]

\subsubsection*{(a) Find $(c)$}

\[
\int_{-\infty}^{\infty} ce^{-2|w|},dw
\]
\[
=====================================
\]
\[
\# 2c\int_0^\infty e^{-2w},dw
\]
\[
\# 2c\cdot \frac{1}{2}
\]
\[
c
\]

Therefore,

\[
c=1
\]

\subsubsection*{(b) Find $(\Pr(-1<W<2))$}

\[
\Pr(-1<W<2)
\]
\[
===========
\]
\[
\int_{-1}^{0} e^{2w},dw
+
\int_{0}^{2} e^{-2w},dw
\]

\[
\# [
\]

\[
\frac{1}{2}(1-e^{-2})+\frac{1}{2}(1-e^{-4})
\]

\[
\# [
\]

\[
1-\frac{1}{2}(e^{-2}+e^{-4})
\]

\subsubsection*{(c) Find $(\Pr(W>0\mid -1<W<2))$}

\[
\Pr(W>0\mid -1<W<2)
\]
\[
===================
\]
\[
\frac{\Pr(0<W<2)}{\Pr(-1<W<2)}
\]

\[
\Pr(0<W<2)
\]
\[
==========
\]
\[
\# \int_0^2 e^{-2w},dw
\]
\[
\frac{1}{2}(1-e^{-4})
\]

Also,

\[
\Pr(-1<W<2)=1-\frac{1}{2}(e^{-2}+e^{-4})
\]

Hence,

\[
\Pr(W>0\mid -1<W<2)
\]
\[
===================
\]
\[
\frac{1-e^{-4}}{2-(e^{-2}+e^{-4})}
\]

---

\section*{9. Laplace Random Variable: Application Examples}

The slides list the following application examples:

\begin{itemize}
\item The modeling of impulsive noise in communication systems, where sudden large noise spikes occur more frequently than in light-tailed models.
\item The error distribution in $(L1)$ (least absolute deviation) regression, where the Laplace distribution arises naturally when minimizing absolute error.
\item The noise added in differential privacy mechanisms, where Laplace noise is used to protect sensitive database information while preserving statistical utility.
\end{itemize}

---

\section*{10. Gamma Random Variable}

The gamma random variable is defined by the following PDF and CDF:

\[
f_X(x)=\frac{(x/b)^{,c-1}\exp(-x/b)}{b\Gamma(c)},u(x)
\]

\[
F_X(x)=\frac{\gamma(c,x/b)}{\Gamma(c)},u(x)
\]

\subsection*{Definitions}

\subsubsection*{Gamma Function}

\[
\Gamma(\alpha)=\int_0^\infty e^{-t}t^{\alpha-1},dt
\]

\subsubsection*{Incomplete Gamma Function}

\[
\gamma(\alpha,\beta)=\int_0^\beta e^{-t}t^{\alpha-1},dt
\]

---

\section*{11. Gamma Random Variable (Graphs and Special Cases)}

The slides show graphs of the gamma PDF and the gamma CDF for multiple parameter choices.

The slides list the following resulting random variables:

\begin{itemize}
\item For $(b=2,\ c=0.5)$, the resulting RV is a Chi-square RV.
\item For $(c=1)$, the resulting RV is an Exponential RV.
\end{itemize}

The slide also states:

\[
\text{where } k=c \text{ and } \theta=b
\]

---

\section*{12. Example \#4 (Gamma RV)}

\subsection*{Problem}

Let $(X)$ be a gamma random variable with parameters $(\alpha)$ and $(\lambda)$. Calculate:

\begin{itemize}
\item (a) $(E[X])$
\item (b) $(Var(X))$
\end{itemize}

\subsection*{Solution}

\subsubsection*{Given}

\[
X\sim Gamma(\alpha,\lambda)
\]

with PDF

\[
f_X(x)=\frac{\lambda^\alpha}{\Gamma(\alpha)}x^{\alpha-1}e^{-\lambda x},
\qquad x\ge 0
\]

\subsubsection*{(a) Find $(E[X])$}

\[
E[X]
\]
\[
====
\]

\[
\# \int_0^\infty x f_X(x),dx
\]

\[
\frac{1}{\Gamma(\alpha)}
\int_0^\infty \lambda x e^{-\lambda x}(\lambda x)^{\alpha-1},dx
\]

\[
\# [
\]

\[
\frac{1}{\lambda \Gamma(\alpha)}
\int_0^\infty \lambda e^{-\lambda x}(\lambda x)^\alpha,dx
\]

\[
\# [
\]

\[
\frac{\Gamma(\alpha+1)}{\lambda \Gamma(\alpha)}
\qquad
(\text{Textbook Eq. }6.1)
\]

\[
E[X]=\frac{\alpha}{\lambda}
\]

\subsubsection*{(b) Find $(Var(X))$}

\[
Var(X)=E[X^2]-(E[X])^2
\]

First compute $(E[X^2])$:

\[
E[X^2]
\]
\[
======
\]

\[
\frac{1}{\Gamma(\alpha)}
\int_0^\infty x^2 \lambda e^{-\lambda x}(\lambda x)^{\alpha-1},dx
\]

\[
\# [
\]

\[
\frac{1}{\lambda^2 \Gamma(\alpha)}
\int_0^\infty \lambda e^{-\lambda x}(\lambda x)^{\alpha+1},dx
\]

\[
\# [
\]

\[
\# \frac{\Gamma(\alpha+2)}{\lambda^2\Gamma(\alpha)}
\]

\[
\frac{\alpha(\alpha+1)}{\lambda^2}
\]

Thus,

\[
Var(X)
\]
\[
======
\]

\[
## \frac{\alpha(\alpha+1)}{\lambda^2}
\]

\[
\# \left(\frac{\alpha}{\lambda}\right)^2
\]

\[
\frac{\alpha}{\lambda^2}
\]

---

\section*{13. Example \#5 (Gamma RV)}

The slide states:

Prove the following properties of the Gamma function.

(a)

\[
\Gamma(n)=(n-1)!
\qquad \text{for } n=1,2,3,\ldots
\]

(b)

\[
\Gamma(x+1)=x\Gamma(x)
\]

(c)

\[
\Gamma(1/2)=\sqrt{\pi}
\]

The slide also states:

Post your answer on Campuswire!

---

\section*{14. Gamma Random Variable: Application Examples}

The slides list the following application examples:

\begin{itemize}
\item The total service time for completing multiple independent tasks, where each task duration is exponentially distributed. The sum of $(k)$ exponential stages follows a Gamma distribution.
\item The received signal energy in detection systems, where the energy computed from multiple independent signal samples follows a Gamma (or Chi-square) distribution.
\item The aggregate waiting time until the $(k)$-th event in a Poisson process, such as the time until the $(k)$-th packet arrival in a network.
\end{itemize}

---

\section*{15. Exercise: Problem Solving 1 (Movie Theater)}

\subsection*{Optimization Problem}

Suppose that you plan to arrive at a movie theater for a scheduled showtime at a fixed start time $(T)$.

\begin{itemize}
\item If you arrive $(s)$ minutes early, you incur a waiting cost of $(c\cdot s)$.
\item If you arrive $(s)$ minutes late, you incur a missing scene cost of $(k\cdot s)$.
\item Suppose also that the travel time from your home to the theater is a continuous random variable with probability density function $(f(t))$.
\end{itemize}

Determine the departure time that minimizes your expected cost.

\subsection*{Solution}

Let $(d)$ denote the departure time from home. Let $(T_r)$ be the random travel time with pdf $(f(t))$ and CDF $(F(t))$.

\subsubsection*{Arrival time}

\[
\text{Arrival}=d+T_r
\]

\subsubsection*{Cost structure}

Early arrival $((d+T_r<T))$: cost $(=c(T-d-T_r))$

Late arrival $((d+T_r>T))$: cost $(=k(d+T_r-T))$

\subsubsection*{Expected cost}

\[
E[C(d)]
\]
\[
=======
\]

\[
\int_0^{T-d} c(T-d-t)f(t),dt
+
\int_{T-d}^{\infty} k(t+d-T)f(t),dt
\]

Differentiate $(E[C(d)])$ with respect to $(d)$ and set to zero:

\[
-cF(T-d)+k(1-F(T-d))=0
\]

\subsubsection*{Optimality condition}

\[
F(T-d^*)=\frac{k}{c+k}
\]

\subsubsection*{Optimal departure time}

\[
d^*=T-F^{-1}!\left(\frac{k}{c+k}\right)
\]

The slide states:

Interpretation: Depart such that the probability of arriving early equals $\left(\dfrac{k}{c+k}\right)$.

---

\section*{16. Exercise: Problem Solving 2 (CDF Analysis)}

The CDF of a RV is given by

\[
F_X(x)=
\begin{cases}
0, & x<0\\
x^2, & 0\le x\le 1\\
1, & x>1
\end{cases}
\]

Find:

\begin{itemize}
\item Mean $((\mu_X))$
\item Variance $((\sigma_X^2))$
\item Skewness $((c_s))$
\item Kurtosis $((c_k))$
\end{itemize}

\subsection*{Solution}

\subsubsection*{Given CDF}

\[
F_X(x)=
\begin{cases}
0, & x<0\\
x^2, & 0\le x\le 1\\
1, & x>1
\end{cases}
\]

\subsubsection*{PDF}

\[
f_X(x)=\frac{d}{dx}F_X(x)=2x,
\qquad 0\le x\le 1
\]

\subsubsection*{Mean}

\[
\mu_X=\int_0^1 x(2x),dx=\frac{2}{3}
\]

\subsubsection*{Variance}

\[
\sigma_X^2
\]
\[
==========
\]

\[
\# \int_0^1 x^2(2x),dx-\mu_X^2
\]

\[
\# \frac{1}{2}-\frac{4}{9}
\]

\[
\frac{1}{18}
\]

\subsubsection*{Skewness}

\[
c_s=\frac{E[(X-\mu)^3]}{\sigma^3}
\]
\[
=================================
\]

\[
-\frac{2\sqrt{2}}{5}
\]

\subsubsection*{Kurtosis}

\[
c_k=\frac{E[(X-\mu)^4]}{\sigma^4}
\]
\[
=================================
\]

\[
\frac{12}{5}
\]

\subsubsection*{Conclusion}

\begin{itemize}
\item Distribution is left-skewed
\item Platykurtic (lighter tails than Gaussian)
\end{itemize}

---

\section*{17. Gaussian Modeling: From Noise to Math}

\subsection*{Motivation}

Why Probability? Because the World is Noisy

\subsection*{The Reality}

\begin{itemize}
\item Sensors don’t give “flat” lines (thermal noise, interference).
\item Repeated measurements form a “cloud” of uncertainty.
\end{itemize}

\subsection*{The Simulation}

\begin{itemize}
\item We model noise using the Gaussian RV: $(X\sim \mathcal{N}(\mu,\sigma^2))$.
\item Generative Formula: $(X=\sigma Z+\mu)$
\end{itemize}

The slide also states:

Coding Task 1: Transform standard randomness $((Z))$ into a physical model $((X))$.

The slide concludes:

Probability models the shape of uncertainty, not just exact values.

---

\section*{18. Density Estimation: Learning from Data}

\subsection*{Reality Check}

We don’t know $(\mu)$ or $(\sigma)$ in the field.

\subsection*{The Problem}

We only observe raw, noisy samples (a histogram).

\subsection*{The Solution}

We use Sample Statistics to estimate the distribution.

The slide presents the following sequence:

\[
\text{Raw Samples } \{x_i\}
\rightarrow
\text{Estimate } \hat{\mu}, \hat{\sigma}
\rightarrow
\text{Estimated PDF}
\]

\subsection*{Key Concept}

We estimate probability distributions, not deterministic functions. This “smooths” the noise into a predictive tool.

---

\section*{19. Applications: Network Systems}

\subsection*{In-Class Focus: Packet Delay Modeling}

\subsection*{Case Study: 500 ICMP Pings (Latency Analysis)}

\subsubsection*{Scenario}

A server transmits packets. Delay varies due to queueing, congestion, and routing hops.

\subsubsection*{Engineering Goal}

\begin{itemize}
\item Typical Delay
\item Jitter (Variation)
\end{itemize}

The slide states:

Critical Discussion: Why is high jitter $((\hat{\sigma}))$ more damaging to real-time video than a high constant delay $((\hat{\mu}))$?

The slide also labels the figure as:

Histogram of observed latency.

---

\section*{20. Applications: Broader Computing Contexts}

\subsection*{Probability in Modern Engineering}

\subsubsection*{Image Processing}

Denoising: Every camera sensor has “grain.” We use Gaussian priors to distinguish between real light and random thermal noise.

Tech: Image Restorers

\subsubsection*{Cloud Systems}

Tail Latency: In distributed systems, we model the “99th percentile” to ensure the slowest $(1\%)$ of users still have a good experience.

Tech: SRE / DevOps

\subsubsection*{IoT \& Robotics}

Sensor Fusion: A GPS is never $(100\%)$ accurate. Kalman filters use probability to “guess” the real location between noisy pings.

Tech: Autonomous Vehicles

The slide concludes:

From hardware sensors to software backends, uncertainty is the only constant.

\end{document}

\end{document}
